\chapter{Chapter Title}

Motivate the section before introducing formal definitions. Explain the analytic role of each
lemma and how the hypotheses correspond to Lean declarations in \texttt{FracLap/}.

\begin{definition}[Descriptive Name]\label{def:example_object}
\lean{FracLap.ExampleObject}
\leanok
Let $s \in (0, 1)$. A function $u \in L^2(\mathbb{R}^d)$ belongs to $\dot{H}^s(\mathbb{R}^d)$ if
its Gagliardo seminorm is finite.
\end{definition}

\begin{lemma}[Intermediate Bound]\label{lem:example_bound}
\lean{FracLap.example_bound}
\uses{def:example_object}
\leanok
Suppose that $u$ satisfies Definition~\ref{def:example_object}. Then there exists a constant
$C = C(d, s) > 0$ such that
\begin{equation}\label{eq:example_bound}
  \lVert u \rVert_{\dot{H}^s(\mathbb{R}^d)} \le C \lVert u \rVert_{H^1(\mathbb{R}^d)}.
\end{equation}
\end{lemma}

\begin{proof}
\uses{def:example_object}
% Add \leanok here only once FracLap.example_bound has zero `sorry`s in Lean.
Splitting the frequency integral at $\lvert \xi \rvert = 1$ gives
\begin{align*}
  \int_{\mathbb{R}^d} \lvert \xi \rvert^{2s} \lvert \hat{u}(\xi) \rvert^2 \dd\xi
  &\le \int_{\lvert \xi \rvert \le 1} \lvert \hat{u}(\xi) \rvert^2 \dd\xi
    + \int_{\lvert \xi \rvert > 1} \lvert \xi \rvert^2 \lvert \hat{u}(\xi) \rvert^2 \dd\xi \\
  &\le \lVert u \rVert_{H^1(\mathbb{R}^d)}^2. \qedhere
\end{align*}
\end{proof}
